\documentclass[11pt]{article}
\usepackage[letterpaper,margin=1in]{geometry}
\usepackage{times}
\usepackage{amsmath}
\usepackage[T1]{fontenc}
\usepackage[hidelinks]{hyperref}
\pagestyle{empty}
\setlength{\parindent}{0pt}
\setlength{\parskip}{6pt}

\begin{document}

\begin{center}
{\Large \bfseries Automated Anatomy-Grounded Lesion Captioning for Whole-Body \\ PSMA PET/CT Using Large Language Models}

\vspace{8pt}
{\large M Mahdi Shakouri$^{1,2}$, Fereshteh Yousefirizi$^{1,2}$}

\vspace{4pt}
{\normalsize $^{1}$ BC Cancer Research Institute, Vancouver, Canada \\
$^{2}$ University of British Columbia, Vancouver, Canada \\
}

\end{center}

\vspace{10pt}
{\bfseries \large Abstract for Technical Review}

Vision language models connecting the text description of an object to its specific location in an
image through visual grounding have shown great potential applications in enhanced radiology reporting.
We present an end-to-end, five-stage pipeline that converts free-text PSMA PET/CT reports and voxel-level lesion segmentations into automatically generated, anatomically grounded, lesion-level captions, together with a multi-metric evaluation framework.

This study was conducted on 286 prostate cancer patient's whole-body $^{68}$Ga/$^{18}$F-PSMA PET/CT scans and corresponding clinical reports. Sentences describing positive findings in PET/CT reports with mentions of standardized uptake values (SUVmax) were identified to form a ground-truth dataset (n=598) of lesion-level captions. 

Manually labelled PET lesions were then querried for corresponding SUVmax values in co-registered PET volumes, and per-lesion SUVmax, three-dimensional centroid, metabolic tumor volume (MTV), and bounding-box dimensions were computed from the volumes. Report-to-lesion correspondence was established via a rank-order matching of SUVmax values within each lesion caption and the patient's PET lesion segmentations. Whole-body CT volumes were segmented into 117 anatomical structures using TotalSegmentator, and for each lesion matched with a caption from the ground-truth dataset, the three nearest anatomical structures were identified via nearest-neighbor search in world coordinates. A structured prompt encoding lesion centroid, SUVmax, MTV, dimensions, and the three nearest anatomic structures (with distances and centroids) was submitted an array of large language models role-framed as a nuclear medicine physician, producing free-text lesion descriptions. 

Generated captions were evaluated against ground-truth report captions using BARTScore, BERTScore, and RaTEScore metrics. In an interim analysis of 384 lesions (from 232 of 286 patients with completed CT segmentation), the pipeline achieved a 79.9\% report-to-lesion correspondence rate and produced captions with a mean BERTScore F1 of 0.872 and RaTEScore of 0.481 against ground-truth report text. This work demonstrates a scalable, reproducible approach to lesion-level caption generation that avoids direct image-token LLM input by grounding generation in structured, automatically derived anatomic and quantitative features.

{\bfseries Keywords:} PSMA PET/CT, prostate cancer, nuclear medicine, lesion captioning, large language models, radiology report generation, automated anatomic segmentation, natural language generation evaluation, multimodal AI, structured reporting

\vspace{6pt}
{\bfseries \large Summary for Program}

We built an automated pipeline that reads PSMA PET/CT lesion data and anatomical context to generate free-text radiology-style descriptions of individual tumor lesions using a large language model, then scored the generated text against the original clinical reports using multiple automated language-quality metrics.

\end{document}
